\chapter{Vis\~ao Geral e Abordagem \`a Solu\c{c}\~ao}
\label{chap:visao-geral}

\noindent{O} presente capítulo descreve a solução proposta nesta dissertação, apresentando a sua arquitetura e os princípios que orientam o seu desenvolvimento. Com base nas limitações identificadas no estado da arte relativamente à aplicação de modelos de linguagem no desenvolvimento de software de forma integrada, define-se uma abordagem estruturada que visa explorar o seu potencial de forma controlada e consistente.

\noindent{A} solução proposta assenta numa pipeline modular, composta por diferentes módulos responsáveis pela geração e transformação de artefactos ao longo do ciclo de vida do desenvolvimento de software. Esta organização permite decompor o processo em etapas distintas, facilitando o controlo, a validação e a evolução independente de cada componente.

\noindent{Para} além da definição da arquitetura global, são também apresentadas as fases que orientaram o desenvolvimento da solução, bem como os principais desafios e decisões de design que influenciaram a sua construção.

\section{Arquitetura e Design da Solu\c{c}\~ao}
\label{sec:visao-arquitetura}

\section{Fases de Desenvolvimento}
\label{sec:visao-fases}

\section{Desafios e Decis\~oes de Implementa\c{c}\~ao}
\label{sec:visao-desafios}
